\section{Thiết kế luồng xử lý và các biểu đồ tuần tự (Sequence Diagrams)}
\label{sec:sequence_diagrams}

Hoạt động phối hợp giữa các thành phần trong hệ thống được chuẩn hóa qua các luồng tương tác tuần tự có kiểm soát trạng thái:

\subsection{Luồng xử lý câu hỏi lâm sàng đa tầng và kiểm soát ranh giới hội thoại}
Quá trình xử lý câu hỏi lâm sàng được phân rã thành chuỗi tương tác tuần tự giữa 6 thành phần: Bác sĩ, Giao diện React, Cổng FastAPI, Bộ kiểm tra an toàn/Truy xuất, Dịch vụ AI và Cơ sở dữ liệu:

\begin{table}[htbp]
\centering
\small
\renewcommand{\arraystretch}{1.15}
\begin{tabularx}{\textwidth}{C{1.0cm} L{3.0cm} L{3.0cm} L{2.5cm} Y}
\toprule
\textbf{Bước} & \textbf{Bên gửi} & \textbf{Bên nhận} & \textbf{Giao thức / Hàm} & \textbf{Mô tả nội dung thông điệp} \\
\midrule
1 & Bác sĩ lâm sàng & Giao diện React & Thao tác UI & Nhập câu hỏi và gửi từ hộp thoại buồng bệnh. \\
\addlinespace
2 & Giao diện React & Cổng FastAPI & \code{POST /api/chat} & Gửi câu hỏi kèm ID ca bệnh và Header xác thực. \\
\addlinespace
3 & Cổng FastAPI & Semantic Cache & Tra cứu SHA-256 & Kiểm tra bộ đệm RAM; nếu trúng đệm, trả về ngay ($< 3\,\text{ms}$). \\
\addlinespace
4 & Cổng FastAPI & NLU Intent Router & Phân loại ý định & Xác định miền câu hỏi; từ chối nếu phát hiện câu hỏi ngoài phạm vi. \\
\addlinespace
5 & NLU Router & INN Normalizer & Chuẩn hóa thực thể & Chuyển biệt dược thương mại sang hoạt chất INN chuẩn hóa. \\
\addlinespace
6 & Cổng FastAPI & Hybrid Retriever & FTS5 + Cosine & Truy xuất đồng thời vector Dược thư trong RAM và từ khóa SQLite. \\
\addlinespace
7 & Hybrid Retriever & Gemini Service & Prompting ($T=0$) & Đóng gói Context y văn vào Prompt, khóa cứng nhiệt độ $0.0$. \\
\addlinespace
8 & Gemini Service & Self-RAG Verifier & Kiểm tra số liệu & Quét tập số liệu trong câu trả lời; nếu có số lạ, kích hoạt ngắt mạch. \\
\addlinespace
9 & Self-RAG Verifier & Giao diện React & Luồng SSE & Phát luồng Server-Sent Events từng token hiển thị giao diện hai vùng. \\
\bottomrule
\end{tabularx}
\caption{Bảng đặc tả tương tác tuần tự luồng xử lý câu hỏi lâm sàng đa tầng}
\label{tab:seq_clinical_query}
\end{table}

\subsection{Luồng kiểm tra tương tác thuốc chống chỉ định tại giao diện CPOE}
Khi bác sĩ chỉ định thuốc trong quá trình kê đơn điện tử:

\begin{table}[htbp]
\centering
\small
\renewcommand{\arraystretch}{1.15}
\begin{tabularx}{\textwidth}{C{1.0cm} L{3.0cm} L{3.0cm} L{2.5cm} Y}
\toprule
\textbf{Bước} & \textbf{Bên gửi} & \textbf{Bên nhận} & \textbf{Giao thức / Hàm} & \textbf{Mô tả nội dung thông điệp} \\
\midrule
1 & Bác sĩ điều trị & Giao diện CPOE & Thao tác UI & Chọn thêm một thuốc mới vào danh sách đơn đang kê. \\
\addlinespace
2 & Giao diện CPOE & Cổng FastAPI & \code{POST /check-ddi} & Gửi danh sách hoạt chất hiện có trong đơn. \\
\addlinespace
3 & Cổng FastAPI & CDSS Sentinel & Tra cứu Hash-Set & Tạo các cặp tổ hợp $\binom{N}{2}$ và tra cứu bảng băm trong RAM ($O(1)$). \\
\addlinespace
4 & CDSS Sentinel & Cổng FastAPI & Kết quả kiểm tra & Trả về danh sách tương tác vi phạm Quyết định 5948/QĐ-BYT (nếu có). \\
\addlinespace
5 & Cổng FastAPI & Giao diện CPOE & HTTP 409 Conflict & Báo lỗi xung đột, gửi kèm danh sách gợi ý thuốc thay thế an toàn. \\
\addlinespace
6 & Giao diện CPOE & Bác sĩ điều trị & Hiển thị cảnh báo & Khóa nút "Lưu đơn", hiển thị thẻ đỏ cảnh báo và nút "Đổi thuốc 1-Click". \\
\addlinespace
7 & Bác sĩ điều trị & Giao diện CPOE & Bấm nút thay thế & Chấp nhận thuốc thay thế; hệ thống tự động kiểm tra lại toàn đơn. \\
\addlinespace
8 & Giao diện CPOE & Bác sĩ điều trị & Trạng thái an toàn & Đơn thuốc hợp lệ; mở khóa trở lại nút "Lưu đơn thuốc". \\
\bottomrule
\end{tabularx}
\caption{Bảng đặc tả tương tác tuần tự luồng kiểm tra tương tác cấm CPOE}
\label{tab:seq_ddi_check}
\end{table}
